% =====================================================================
\ProjectChapter{Installing and Running the Languages}{app:install}
\setcounter{table}{0}\renewcommand{\thetable}{R\arabic{table}}
% =====================================================================

All languages in these notes except MATLAB and the full Mathematica can be installed for
free on Windows, macOS and Linux. Table~\ref{tab:install} lists the recommended tools. On
Windows, the \textbf{Windows Subsystem for Linux} (WSL) gives you a complete Ubuntu Linux
environment in which every Linux command in these notes works unchanged; install it with
\code{wsl --install} in an administrator PowerShell. On macOS, the package manager
\textbf{Homebrew} (\url{https://brew.sh}) installs most tools with one command.

\begin{table}[htbp]
\centering
\caption{Getting each language. ``Online'' lists a way to try the language in a web browser
without installing anything.}
\label{tab:install}
\small
\begin{tabularx}{\linewidth}{@{}l>{\RaggedRight\arraybackslash}X>{\RaggedRight\arraybackslash\ttfamily}p{3.9cm}>{\RaggedRight\arraybackslash}p{3.0cm}@{}}
\toprule
\textbf{Language} & \textbf{Install} & \normalfont\textbf{Ubuntu / WSL} & \textbf{Online}\\
\midrule
\py & Anaconda / Miniforge, python.org, uv & sudo apt install python3-pip & Google Colab, JupyterLite\\
Fortran & gfortran (GCC), Intel oneAPI ifx, LFortran & sudo apt install gfortran & godbolt.org, play.fortran-lang.org\\
C, C++ & GCC, Clang, Visual Studio (Windows), Xcode tools (macOS) & sudo apt install\newline build-essential & godbolt.org\\
C\# & .NET SDK (dotnet.microsoft.com) & sudo apt install dotnet-sdk-8.0 & dotnetfiddle.net\\
Java & Eclipse Temurin JDK (adoptium.net) & sudo apt install default-jdk & JDoodle, OneCompiler\\
JavaScript & any web browser; Node.js (nodejs.org) & sudo apt install nodejs & browser console\\
MATLAB & university licence (mathworks.com) & {\normalfont n/a} & MATLAB Online\\
Octave & octave.org & sudo apt install octave & octave-online.net\\
Wolfram & Mathematica (licence) or free Wolfram Engine & {\normalfont n/a} & Wolfram Cloud, Wolfram|Alpha\\
R & cran.r-project.org + RStudio/Positron & sudo apt install r-base & posit.cloud, webR\\
Julia & juliaup (julialang.org) & see julialang.org & JuliaHub\\
Rust & rustup (rustup.rs) & see rustup.rs & play.rust-lang.org\\
Go & go.dev/dl & sudo apt install golang & go.dev/play\\
SQL (SQLite) & sqlite.org; built into \py & sudo apt install sqlite3 & sqliteonline.com\\
Modelica & OpenModelica (openmodelica.org) & see website & n/a\\
\bottomrule
\end{tabularx}
\end{table}

\section*{One editor for everything}

\textbf{Visual Studio Code} (\url{https://code.visualstudio.com}) is a free editor that
supports every language in these notes through extensions: syntax highlighting, code
completion, running and debugging, Jupyter notebooks, Git integration and remote work on
WSL or a cluster. Install the extensions for \py, C/C++, Fortran (``Modern Fortran''), C\#,
Java, Julia, R, Rust (``rust-analyzer'') and Go as you need them.

\section*{Running the example programs}

All example programs of these notes are in the folder \file{code/} of the repository,
one sub-folder per language, together with the data files in \file{code/data/}. The
commands used to compile and run the drag program in each language were:
\begin{lstlisting}[style=shell]
$ python3 drag.py                                 # Python
$ gfortran -O2 -Wall -o drag drag.f90 && ./drag   # Fortran
$ gcc -std=c11 -O2 -Wall -o drag drag.c -lm && ./drag     # C
$ g++ -std=c++20 -O2 -Wall -o drag drag.cpp && ./drag     # C++
$ mcs Drag.cs && mono Drag.exe                    # C# with Mono
$ dotnet run Drag.cs                              # C# with .NET 10 or later
$ java Drag.java                                  # Java 11 or later
$ node drag.js                                    # JavaScript
$ octave --eval drag                              # Octave; in MATLAB type: drag
$ Rscript drag.R                                  # R
$ julia drag.jl                                   # Julia
$ rustc -O drag.rs && ./drag                      # Rust
$ go run drag.go                                  # Go
$ wolframscript -file drag.wl                     # Wolfram Language
\end{lstlisting}
The programs that read a data file look for it in the folder \file{code/data/}. In the
programs this folder is written as \file{../data/}, so run each program from inside its
own language folder. The programs of Chapter~\ref{ch:rosetta} are the exception: they read
\file{yield.txt} from the current folder, and a copy of that file is kept in
\file{code/rosetta/} for this purpose.

To re-run every example at once and compare each output with the one printed in these
notes, run the checking script from the top folder of the repository. Languages that are
not installed on your computer are skipped.
\begin{lstlisting}[style=shell]
$ python3 tools/check_examples.py
\end{lstlisting}

% =====================================================================
\ProjectChapter{Python Quick Reference}{app:pyref}
% =====================================================================

\section*{Basics}
\begin{lstlisting}[language=Python]
x = 3.0; n = 10; name = "steel"; ok = True; nothing = None
type(x); int("3"); float("2.5"); str(42); round(3.14159, 2)
7 / 2   # 3.5      7 // 2  # 3      7 % 2  # 1      2 ** 10  # 1024
a == b; a != b; a < b <= c; x and y; x or y; not x
print(f"{x:8.3f} {n:5d} {x:.2e} {name:>10s}")
value = input("Prompt: ")          # always a string
\end{lstlisting}

\section*{Control flow}
\begin{lstlisting}[language=Python]
if x > 0:
    ...
elif x == 0:
    ...
else:
    ...
for i in range(n):          # 0 ... n-1
for i in range(1, n + 1):   # 1 ... n
for item in items:
for i, item in enumerate(items):
for a, b in zip(list_a, list_b):
while condition:
break; continue; pass
result = "big" if x > 10 else "small"      # conditional expression
\end{lstlisting}

\section*{Collections}
\begin{lstlisting}[language=Python]
L = [1, 2, 3]; L.append(4); L.insert(0, 0); L.pop(); L.sort(); len(L)
L[0]; L[-1]; L[1:3]; L[::-1]; 3 in L
squares = [x**2 for x in range(10) if x % 2 == 0]
t = (1, 2); a, b = t
d = {"steel": 7850}; d["al"] = 2700; d.get("ti", 0); d.keys(); d.items()
s = {1, 2, 3}; s.add(4); s & {2, 5}; s | {9}
\end{lstlisting}

\section*{Functions, classes, errors, files}
\begin{lstlisting}[language=Python]
def f(x, y=1.0, *args, **kwargs):
    """Docstring."""
    return x + y
g = lambda x: x**2

class Tank:
    def __init__(self, volume):
        self.volume = volume
    def mass(self, rho):
        return rho * self.volume

try:
    risky()
except (ValueError, ZeroDivisionError) as err:
    print("error:", err)
finally:
    cleanup()
raise ValueError("message"); assert x > 0, "x must be positive"

with open("data.txt") as f:            # "w" write, "a" append
    for line in f:
        values = [float(v) for v in line.split()]
\end{lstlisting}

\section*{NumPy, Matplotlib and SciPy}
\begin{lstlisting}[language=Python]
import numpy as np, matplotlib.pyplot as plt
a = np.array([1., 2., 3.]); np.zeros((3, 3)); np.ones(5); np.eye(3)
np.arange(0, 1, 0.1); np.linspace(0, 1, 11); np.random.default_rng(0).normal(size=5)
a.shape; a.reshape(3, 1); a.T; a[a > 1]; A[1, :]; A[:, 0]
a.sum(); a.mean(); a.std(ddof=1); a.max(); np.argmax(a); np.cumsum(a)
np.sin(a); np.exp(a); np.sqrt(a); np.abs(a); np.isclose(a, b)
A @ B; np.linalg.solve(A, b); np.linalg.eig(A); np.linalg.norm(a)
np.loadtxt("f.txt"); np.savetxt("out.txt", a); np.genfromtxt("f.csv", delimiter=",")

fig, ax = plt.subplots()
ax.plot(x, y, "o-", label="data"); ax.set_xlabel("x (m)"); ax.legend()
fig.savefig("plot.pdf")

from scipy.optimize import brentq, fsolve, minimize, curve_fit
from scipy.integrate import quad, solve_ivp
from scipy.interpolate import interp1d
\end{lstlisting}

% =====================================================================
\ProjectChapter{Numerical Functions Across Languages}{app:functions}
% =====================================================================

\begin{sidewaystable}
\centering
\caption{Common mathematical and numerical operations.}
\small
\setlength{\tabcolsep}{3pt}
\resizebox{\linewidth}{!}{%
\begin{tabular}{@{}>{\bfseries}l>{\ttfamily}l>{\ttfamily}l>{\ttfamily}l>{\ttfamily}l>{\ttfamily}l@{}}
\toprule
 & \normalfont\bfseries \py & \normalfont\bfseries MATLAB & \normalfont\bfseries R & \normalfont\bfseries Julia & \normalfont\bfseries C / C++\\
\midrule
$\sqrt{x}$ & math.sqrt(x) & sqrt(x) & sqrt(x) & sqrt(x) & sqrt(x)\\
$e^x$, $\ln x$ & math.exp, math.log & exp, log & exp, log & exp, log & exp, log\\
$\log_{10} x$ & math.log10(x) & log10(x) & log10(x) & log10(x) & log10(x)\\
$|x|$ & abs(x) & abs(x) & abs(x) & abs(x) & fabs(x) / std::abs\\
$\pi$ & math.pi & pi & pi & pi, \textbackslash pi & M\_PI / std::numbers::pi\\
$\operatorname{atan2}(y,x)$ & math.atan2(y, x) & atan2(y, x) & atan2(y, x) & atan(y, x) & atan2(y, x)\\
round & round(x) & round(x) & round(x) & round(x) & round(x)\\
remainder & x \% y & mod(x, y) & x \%\% y & mod(x, y) & fmod(x, y)\\
sum & np.sum(a) & sum(a) & sum(a) & sum(a) & std::accumulate\\
mean, std & np.mean, np.std(ddof=1) & mean, std & mean, sd & mean, std & (by hand)\\
min, max & np.min, np.max & min, max & min, max & minimum, maximum & std::min\_element\\
sort & np.sort(a) & sort(a) & sort(a) & sort(a) & std::sort\\
cumulative sum & np.cumsum(a) & cumsum(a) & cumsum(a) & cumsum(a) & std::partial\_sum\\
dot product & a @ b & dot(a, b) & sum(a*b) & dot(a, b) & std::inner\_product\\
matrix product & A @ B & A*B & A \%*\% B & A*B & (Eigen) A*B\\
transpose & A.T & A\textquotesingle{} & t(A) & A\textquotesingle{} & (Eigen) A.transpose()\\
solve $A\mathbf{x}=\mathbf{b}$ & np.linalg.solve & A\textbackslash b & solve(A, b) & A\textbackslash b & (LAPACK / Eigen)\\
eigenvalues & np.linalg.eigvals & eig(A) & eigen(A) & eigvals(A) & (LAPACK / Eigen)\\
root of $f$ & scipy.optimize.brentq & fzero & uniroot & (Roots.jl) find\_zero & (GSL)\\
integral & scipy.integrate.quad & integral & integrate & (QuadGK.jl) quadgk & (GSL)\\
ODE & solve\_ivp & ode45 & deSolve::ode & (DiffEq.jl) solve & (Boost.Odeint)\\
interpolate & np.interp & interp1 & approx & (Interpolations.jl) & {\normalfont n/a}\\
polynomial fit & np.polyfit & polyfit & lm(y \textasciitilde\ poly(x, n)) & (Polynomials.jl) fit & {\normalfont n/a}\\
random uniform & rng.random() & rand & runif(1) & rand() & std::uniform\_real\_distribution\\
random normal & rng.normal() & randn & rnorm(1) & randn() & std::normal\_distribution\\
read table & np.loadtxt, pd.read\_csv & readmatrix, readtable & read.csv & CSV.read & (by hand)\\
\bottomrule
\end{tabular}}
\end{sidewaystable}

\noindent In C and C++, include \code{<math.h>}/\code{<cmath>} for the mathematical functions.
Fortran uses the same names as C for most intrinsics (\code{sqrt}, \code{exp}, \code{log},
\code{log10}, \code{abs}, \code{atan2}, \code{mod}) plus array intrinsics (\code{sum},
\code{maxval}, \code{matmul}, \code{dot\_product}, \code{transpose}).

% =====================================================================
\ProjectChapter{Glossary}{app:glossary}
% =====================================================================

\begin{description}[style=nextline,leftmargin=1.2em,font=\sffamily\bfseries]
\item[Algorithm] A finite, precise sequence of steps that solves a problem.
\item[API] Application programming interface: the set of functions, classes and
  conventions through which one piece of software is used by another.
\item[Array] An ordered collection of elements of the same type, accessed by index.
\item[Assembly language] A human-readable form of machine code, specific to one processor
  family.
\item[Broadcasting] Combining arrays of different shapes by automatically extending the
  smaller one (NumPy, Julia, MATLAB).
\item[Bug] An error in a program. \textbf{Debugging} is finding and fixing it.
\item[Bytecode] An intermediate, platform-independent code executed by a virtual machine
  or interpreter.
\item[Class] A blueprint for objects, defining their data and methods.
\item[Column-major] Storing a matrix column by column in memory (Fortran, MATLAB, R,
  Julia).
\item[Compiler] A program that translates source code into machine code (or bytecode)
  before the program runs.
\item[Debugger] A tool that lets you pause a running program, step through it and inspect
  variables.
\item[Dynamic typing] Types belong to values and are checked while the program runs.
\item[Exception] An object that signals a run-time error and can be caught and handled.
\item[Floating-point number] A computer representation of a real number with a limited
  number of significant digits (IEEE 754).
\item[Function] A named, reusable block of code with inputs (parameters) and usually an
  output (return value). Also called a procedure, subroutine or method.
\item[Garbage collector] Part of a language runtime that automatically frees memory that is
  no longer used.
\item[Heap] The memory region for data whose size or lifetime is decided at run time.
\item[IDE] Integrated development environment: editor, runner, debugger and other tools in
  one program.
\item[Immutable] Cannot be changed after creation (e.g.\ \py strings and tuples).
\item[Inheritance] Defining a new class as an extension of an existing one.
\item[Interpreter] A program that executes source code (or bytecode) directly, statement
  by statement.
\item[JIT compilation] Just-in-time compilation: translating code to machine code while
  the program runs, just before it is needed.
\item[Library] A collection of reusable, tested code (functions, classes) for a task.
\item[Linker] A program that combines compiled object files and libraries into an
  executable.
\item[Literal] A value written directly in the code, such as \code{3.14} or
  \code{"steel"}.
\item[Machine epsilon] The gap between 1 and the next larger floating-point number;
  $\approx 2.2\times10^{-16}$ in double precision.
\item[Memory leak] Memory that a program allocates but never frees.
\item[Method] A function that belongs to a class or object.
\item[Module] A file or unit of code that groups related definitions and can be imported.
\item[Multiple dispatch] Choosing which version of a function to run based on the types of
  all its arguments (Julia).
\item[Notebook] A document mixing executable code cells, results and text (Jupyter,
  Mathematica, MATLAB Live Scripts, Pluto).
\item[Object] A bundle of data and the methods that act on it; an instance of a class.
\item[Off-by-one error] A loop or index that is one step too short or too long.
\item[Operator overloading] Defining what operators such as \code{+} and \code{*} mean for
  user-defined types.
\item[Overflow] A result too large for the type that must hold it.
\item[Package manager] A tool that installs libraries and their dependencies (\code{pip},
  \code{conda}, CRAN, \code{npm}, Cargo, NuGet, Maven, Julia's Pkg).
\item[Paradigm] A style of programming: imperative, procedural, object-oriented,
  functional, declarative, array.
\item[Pointer] A variable holding a memory address (C, C++, Fortran pointers).
\item[Polymorphism] Code that works with objects of different types through a common
  interface.
\item[Profiling] Measuring where a program spends its time.
\item[Recursion] A function that calls itself.
\item[Reference] A name that refers to an object without copying it.
\item[REPL] Read--evaluate--print loop: an interactive prompt.
\item[Row-major] Storing a matrix row by row in memory (C, C++, \py/NumPy by default).
\item[Runtime] The period when a program runs; also the support software a language needs
  while running (JVM, .NET runtime).
\item[Scope] The region of a program in which a name is visible.
\item[Source code] The human-readable text of a program.
\item[Stack] The memory region holding the local variables of active function calls.
\item[Standard library] The library that comes with a language implementation.
\item[Static typing] Types are attached to variables and checked before the program runs.
\item[String] A sequence of characters (text).
\item[Syntax] The grammar rules of a language.
\item[Type] The kind of a value (integer, float, string, \ldots), which determines how it is
  stored and what can be done with it.
\item[Undefined behaviour] In C and C++, an operation whose result the language standard
  does not define (e.g.\ reading outside an array); anything may happen.
\item[Unicode, UTF-8] The universal character set and its most common byte encoding.
\item[Variable] A name that refers to a value stored in memory.
\item[Vectorisation] Expressing a computation as operations on whole arrays instead of
  explicit loops; also the use of SIMD processor instructions.
\item[Version control] Recording the history of changes to files (Git).
\item[Virtual machine] Software that executes bytecode (JVM, .NET CLR).
\end{description}

